\subsection{全子代数}

设 A 是 K 上的一个幺代数。

定义 3. --- A 的全子代数是指 A 的一个幺子代数 B，使得 B 中每个在 A 内可逆的元素在 B 内也可逆。

换言之，B 是 A 的全子代数，当且仅当 \(\mathrm{Sp}_{\mathrm{B}}(x)=\mathrm{Sp}_{\mathrm{A}}(x)\) 对一切 \(x\in B\) 成立。

A 的一族全子代数的交是 A 的一个全子代数。

设 M 是 A 的一个子集。A 的包含 M 的诸全子代数之交，是 A 的包含 M 的最小全子代数。

它称为由 M 生成的 A 的全子代数。M 在 A 中的交换子 M′ 是 A 的一个全子代数（因为，若 x 在 A 中可逆且与 M 交换，则 \(x^{-1}\) 与 M 交换）。因此，M 的双交换子 M'' 是 A 的一个全子代数，它包含由 M 生成的 A 的全子代数。

若 M 的元素两两可交换，我们有 \(M \subset M'\) 与 \(M'' \subset M'''\)；于是代数 \(M''\) 是交换的，从而由 M 生成的全子代数也是交换的。

A 的一个极大交换子代数是全子代数，因为它等于其交换子。

引理 2. --- 设 \((x_{\lambda})_{\lambda \in \Lambda}\) 是 A 中两两交换的一族元素。由 \((x_{\lambda})\) 生成的全子代数是由形如 \(\mathrm{R}((x_{\lambda}))\) 的元素所成的集合，其中 \(\mathrm{R} \in \mathrm{K}((\mathrm{X}_{\lambda}))\) 取遍使族 \((x_{\lambda})\) 可代入的那些有理分式。

设 B 是由族 \((x_{\lambda})\) 生成的 A 的全子代数，并用 \(B_{1}\) 表示由形如 \(\mathrm{R}((x_{\lambda}))\) 的元素所成的集合，其中 \(\mathrm{R}((\mathrm{X}_{\lambda}))\) 是使 \((x_{\lambda})\) 可代入的有理分式。明确地说，\(B_{1}\) 是 A 中形如 \(\mathrm{P}((x_{\lambda}))\mathrm{Q}((x_{\lambda}))^{-1}\) 的元素所成的集合，其中 \(\mathrm{P},\mathrm{Q}\in\mathrm{K}[(\mathrm{X}_{\lambda})]\) 且 \(\mathrm{Q}((x_{\lambda}))\) 在 A 中可逆。集合 \(B_{1}\) 是 A 的包含族 \((x_{\lambda})\) 的一个幺子代数。它是一个全子代数：若 \(\mathrm{P}((x_{\lambda}))\mathrm{Q}((x_{\lambda}))^{-1}\) 在 A 中可逆，则 \(\mathrm{P}((x_{\lambda}))\) 在 A 中可逆，且逆 \(\mathrm{Q}((x_{\lambda}))\mathrm{P}((x_{\lambda}))^{-1}\)（即 \(\mathrm{P}((x_{\lambda}))\mathrm{Q}((x_{\lambda}))^{-1}\) 的逆元）属于 \(B_{1}\)。于是我们有 \(B\subset B_{1}\)。另一方面，若 \(P,Q\in K[(X_{\lambda})]\)，且 \(\mathrm{Q}((x_{\lambda}))\) 在 A 中可逆，则 \(\mathrm{P}((x_{\lambda}))\in\mathrm{B}\) 且 \(\mathrm{Q}((x_{\lambda}))\in\mathrm{B}\)，因而 \(\mathrm{Q}((x_{\lambda}))^{-1}\in\mathrm{B}\) 且 \(\mathrm{P}((x_{\lambda}))\mathrm{Q}((x_{\lambda}))^{-1}\in\mathrm{B}\)，从而 \(B_{1}\subset B\)。

